\color{unidarkgreen}{\subsection{Логика управления топологией}}
\color{black}

\begin{enumerate}[label=\arabic{section}.\arabic{subsection}.\arabic{enumi}, labelsep=4pt, leftmargin=0pt, itemindent=57pt, itemsep=0pt, parsep=5pt] % если что можно крутить itemsep


%%%%%%%%%%%%%%%%%%%%%%%%%%%%%%%%%%%%%%%%%%%%%%%%%% Логика управления топологией %%%%%%%%%%%%%%%%%%%%%%%%%%%%%%%%%%%%%%%%%%%%%%%%%%
\item Общие данные

\begin{enumerate}[label=\arabic{section}.\arabic{subsection}.\arabic{enumi}.\arabic*, labelsep=4pt, leftmargin=0em, itemindent=65pt, parsep=0pt]

\item
Логика управления топологией предназначена для назначения фиксации присоединений к СШ в соответствии со схемой первичных соединений.
С целью учета различных особенностей выполнения схем СШ в устройствах все присоединения имеют одинаковые возможности в задании конфигурации, т.е. каждое присоединение может быть зафиксировано за первой или второй СШ или назначено в качестве ШСВ посредством программного переключателя \textbf{«Фиксация Вn»}, где n – номер присоединения:
\begin{itemize} 
\item состояние \textbf{«Внеш.»} применяется для фиксации присоединения по состоянию двух входных сигналов: «Вn к 1С» и «Вn к 2С». В качестве внешних сигналов, например, могут выступать положения блок-контактов разъединителей на присоединении.
При пропадании срабатывания внешних сигналов логика сохраняет последнюю фиксацию присоединения. Сброс фиксации происходит при отсутствии состояния \textbf{«Внеш.»} (переход на другое состояние) или при срабатывании смежного внешнего сигнала.
\item состояние \textbf{«Вывод»} применяется для вывода присоединения из схемы ДЗШ. В этом состоянии ток присоединения не участвует в расчете дифференциальных и тормозных токов, используемых дифференциальными органами.
При выводе выключателя присоединения в ремонт (при поступлении внешнего сигнала \textit{«Ремонт Вn»}) ток присоединения также исключается из расчетов дифференциальных и тормозных токов.
\item состояние \textbf{«1С»} применяется для жесткой фиксации присоединения за первой СШ. В этом состоянии ток присоединения участвует в расчете дифференциальных и тормозных токов ИО1 и ПО.
\item состояние \textbf{«2С»} применяется для жесткой фиксации присоединения за второй СШ. В этом состоянии ток присоединения участвует в расчете дифференциальных и тормозных токов ИО2 и ПО.
\item состояние \textbf{«ШСВ1»} устанавливается для присоединения типа ШСВ в схеме ДЗШ с двумя ТТ ШСВ. В это состоянии фиксируется ТТ, установленный со стороны второй СШ. Присоединение учитывается за первой СШ, и его ток включается в расчет дифференциального и тормозного токов ИО1.
\item состояние \textbf{«ШСВ2»} устанавливается для присоединения типа ШСВ в схеме ДЗШ с двумя ТТ ШСВ. В это состоянии фиксируется ТТ, установленный со стороны первой СШ. Присоединение учитывается за второй СШ, и его ток включается в расчет дифференциального и тормозного токов ИО2.
\item состояние \textbf{«ШСВ12»} используется в схемах ДЗШ с одним ТТ ШСВ, установленного со стороны второй СШ. В этом состоянии присоединение учитывается за первой и второй СШ: ток включается в расчет дифференциальных и тормозных токов ИО1 со знаком «плюс», а в расчет токов ИО2 – со знаком «минус».
\item состояние \textbf{«ШСВ21»} используется в схемах ДЗШ с одним ТТ ШСВ, установленного со стороны первой СШ. В этом состоянии присоединение учитывается за первой и второй СШ: ток включается в расчет дифференциальных и тормозных токов ИО2 со знаком «плюс», а в расчет токов ИО1 – со знаком «минус».
\end{itemize}

\medskip

\begin{figure}[H]
\centering
\includegraphics[width=0.9\textwidth,height=0.9\textheight,keepaspectratio]{\fbpath/661.1721 ЛУТ/img/Fix_Q1.pdf}
\captionof{figure}{Функциональный блок органа фиксации В1}\label{pict:Fix_Q1}
\end{figure}

\medskip

\item
Режим нарушенной фиксации инициируется либо оперативно – посредством внешнего сигнала \textit{«Нарушение фикс.»}, либо автоматически – через сигналы, указывающие на нарушение фиксации конкретного присоединения.

\item
Сигнал нарушения фиксации отдельного присоединения может быть сформирован в следующих случаях:
\begin{itemize}
\item при отсутствии сигнала \textit{«Ремонт Вn»}, когда выключатель присоединения не находится в ремонте, фиксация осуществляется от внешних сигналов (положение \textbf{«Внеш.»}), и срабатывают оба сигнала фиксации присоединений: \textit{«Вn к 1С»} и \textit{«Вn к 2С»};
\item при фиксации присоединений по внешним сигналам (положение \textbf{«Внеш.»}), если в данном состоянии внешние сигналы фиксации исчезнут, но по присоединению продолжит протекать ток.
\end{itemize}

\item
Выходной сигнал \textit{«Нарушение фикс.»} удерживается в течение 200 мс после восстановление фиксации.
Данная выдержка необходима для исключения влияния переходных состояний, возникающих при переключении разъединителей или кратковременном пропадании входного сигнала \textit{«Вn к 1С»} (\textit{«Вn к 2С»}).

\medskip

\begin{figure}[H]
\centering
\includegraphics[width=0.9\textwidth,height=0.9\textheight,keepaspectratio]{\fbpath/661.1721 ЛУТ/img/LUT.pdf}
\captionof{figure}{Функциональный блок логики нарушение фиксации}\label{pict:LUT}
\end{figure}

\medskip

\item
Уставки функции фиксации присоединений представлены в таблице \ref{desc:tbl1}.

\small
\begin{longtable}{|>{\centering\arraybackslash}m{5.3cm}|>{\centering\arraybackslash}m{3.3cm}|>{\centering\arraybackslash}m{4.2cm}|>{\centering\arraybackslash}m{1.8cm}|>{\centering\arraybackslash}m{1cm}|}
% Первая страница: полный заголовок
\caption{Перечень уставок функции фиксации присоединений\hfill\vspace{-0.5\baselineskip}}\label{desc:tbl1}\\
\hline
\rowcolor{gray!30}
Параметр & Условное обозначение на схеме & Значение/ Диапазон & Единица измерения & Шаг \\
\hline
\endfirsthead

% Промежуточные страницы: «Продолжение таблицы»
\caption*{\hspace{3pt}\emph{Продолжение таблицы \ref{desc_dzl:tbl1}}\hfill\ }\vspace{-0.5\baselineskip} \\
\hline
\rowcolor{gray!30}
Параметр & Условное обозначение на схеме & Значение/ Диапазон & Единица измерения & Шаг \\
\hline
\endhead

%===m>ДЗШ|-
%===m
\end{longtable}
\normalsize	


\end{enumerate}

\end{enumerate}
